\section{High Treewidth, Low SDP-width}\label{sec:HighTlowS}

In this section, we introduce the notion of \emph{symbolic dynamic programming
width} (SDP-width) of a graph with respect to a DP-core $\dpcore$. Let $\dpcore$
be a DP-core solving a graph problem $\problem$, and let $\treedec$ be a tree
decomposition of a graph $\graphG$. We define the SDP-width
$\sw(\dpcore,\treedec)$ to be the maximum width of a (minimized) OBDD representing
a set of partial solutions at any node of $\treedec$.
We define the SDP-width of a graph $\graphG$ with respect to $\dpcore$ as the
minimum SDP-width over all tree decompositions of $\graphG$:

$$\sw(\dpcore,\graphG)
= \argmin_{\treedec}\sw(\dpcore,\treedec).$$


We note that SDP-width depends on the DP-core $\dpcore$, particular in the
chosen (binary) encoding of witnesses representing partial solutions for the
problem $\problem$. It is also sensitive to the choice of tree decomposition
$\treedec$, as this determines the subgraphs induced at the bags of
$\treedec$. For example, a $k$-instructive path decomposition of an $n$-vertex
complete graph that introduces vertices and edges in an arbitrary order may
contain a bag whose induced subgraph is an arbitrary $n$-vertex graph, including
graphs with unbounded SDP-width. This motivates the study of graph classes for which high treewidth does not
necessarily imply large SDP-width under suitably chosen decompositions.

A tree decomposition is \emph{immediate} if whenever a node $u$ is an introduce node, introducing a vertex $v$, all edges incident with $v$ in the graph $G_u$ are subsequently introduced and vertex $v$ is forgotten as soon all edges of $\graphG$ incident to $\vertexv$ are introduced. 
% \section{High Treewidth, Low SDP-width}\label{sec:HighTlowS}

% In this section we introduce the notion of \emph{symbolic dynamic programming width} (SDP-width) of a graph with respect to a Dp-core $\dpcore$. Given a DP-core $\dpcore$ solving graph problem $\problem$ and tree decomposition $\treedec$ of $\graphG$, we let the SDP-width, $\sw(\dpcore, \treedec)$, be the width of the highest width (minimized) OBDD representing a set of partial solutions at a node of $\treedec$. In turn, we define the SDP-width of graph $\graphG$ to be the minimum SDP-width over all tree decompositions over $\graphG$, 
% $$\sw(\dpcore, \graphG) = \argmin_{\treedec}\sw(\dpcore, \treedec).$$ 
% We note that SDP-width depends on DP-core $\dpcore$ particularly in the chosen (binary) encoding of witnesses comprising sets of partial solutions for problem $\problem$. SDP-width is also sensitive to tree decomposition $\treedec$ particularly in the subgraphs induced at bags of $\treedec$. For instance, a $k$-instructive path decomposition of an $n$-node complete graph that arbitrarily introduces nodes and edges may include a bag representing an $n$-node graph from any graph class including whatever/any graph class realizing unbounded SDP-width.

In the remainder of this section, we survey classes of graphs that exhibit high treewidth while either analytically or  empirically (see section \ref{sec:Experiments}) admitting compact symbolic representations over immediate tree decompositions.

\paragraph*{Complete graphs}
For each $n \in \N$, the complete graph on $n$ vertices, denoted $K_n$, is the graph $\graphG = (\vertexset,\edgeset)$ where $\vertexset = [n]$ and $\edgeset = \binom{[n]}{2}$. That is, $K_n$ contains an edge between every pair of distinct vertices. It is well known that the treewidth of $K_n$ is $n-1$.

\begin{observation}[Folklore]
    The \emph{complete graph} on $n$ vertices has treewidth $n-1$.
\end{observation}

For the case of complete graphs the we are able to bound SDP-width analytically. 

\begin{proposition} \todo{Create appendix section for proofs in section \ref{sec:HighTlowS}.}
\label{proposition:WidthClique}
Let $\dpcore$ be a DP-core solving {\sc Vertex Cover} and let $\treedec$ be an immediate tree decomposition of complete graph $K_n$. Then $\sw(\dpcore, \treedec)$ is bounded above by $\BigO(1)$.\todo{Should I write big theta here? Its obvious that I will never do better than $\BigO(1)$ isn't it?}
\end{proposition}

In particular the set of partial solutions for {\sc Vertex Cover} at any node of an immediate path decomposition can be represented by an OBDD of width 4. We note that Proposition \ref{proposition:WidthClique} holds over any immediate tree decomposition and is independent of variable ordering.

 
\paragraph*{Grid graphs}
For $r,s \in \N^+$, the \emph{$r \times s$ grid graph}, denoted $G_{r,s}$, is the graph $\graphG = (\vertexset,\edgeset)$ where $\vertexset = [r] \times [s]$, and $\edgeset = \bigl\{\{(i,j),(i',j')\} \mid |i-i'| + |j-j'| = 1 \bigr\}$. Intuitively, the vertices of $G_{r,s}$ can be arranged in an $r \times s$ grid, where each vertex $(i,j)$ is adjacent to its horizontally and vertically
neighboring grid points. Grid graphs have treewidth $\Theta(\min\{r,s\})$. 

\begin{observation}[Folklore]
    The $r \times s$ grid graph has treewidth $\Theta(\min\{r,s\})$.
\end{observation}
In particular the family of square grids $G_{n,n}$ has treewidth $\Theta(n)$.


\paragraph*{Rook graphs}
For $r \in \N^+$, the \emph{rook graph} $R_r$ is the graph
$\graphG = (\vertexset,\edgeset)$ where $\vertexset = [r] \times [r]$, and $\edgeset = \bigl\{\{(i,j),(i',j')\} \mid i=i' \text{ or } j=j',\ (i,j) \neq (i',j') \bigr\}$.
Intuitively, the vertices can be arranged in an $r \times r$ grid, where any two vertices in the same row or the same column are adjacent, so that each row and each column induces a clique. Rook graph $R_r$ is equivalently characterized as the line graph of the complete bipartite graph $K_{r,r}$ and or the 2 dimensional Hamming graph over set  $[r]$. Under each of these characterizations, $R_r$ rook graphs have been shown to have treewidth $\Theta(r^2)$ (See Theorem 1.2 in \cite{HarveyTreewidth2014} and Theorem 3 in \cite{SUNILCHANDRAN} respectively).

\begin{observation}
    The $N$-node rook graph $R_n$ (with $N = n^2$) has treewidth $\Theta(N)$.
\end{observation}

The SDP-width of can be bounded analytically. In particular, given an $N$-node rook graph the set of partial solutions for {\sc Vertex Cover} at any node of an immediate path decomposition can be represented by an OBDD of width at most $2^{\BigO(\sqrt{N})}$.

\begin{proposition}
\label{proposition:SDPWidthRook}
Let $\dpcore$ be a DP-core solving {\sc Vertex Cover} and let $\treedec$ be an immediate tree decomposition of the $N$-node rook graph $R_n$ (with $N = n^2$). Then $\sw(\dpcore, \treedec)$ is bounded above by $2^{\BigO(\sqrt{N})}$.
\end{proposition}

We note that the number of optimal solutions to {\sc Vertex Cover} on $N$-node rook graphs is given by $(\sqrt{N})! \gg 2^{\sqrt{N}}$. We note that Proposition \ref{proposition:SDPWidthRook} holds over any immediate tree decomposition and is independent of variable ordering.



% \begin{proposition}
% \label{proposition:WidthClique}
% Let $K_n$ be the clique graph with $n$ vertices. The final set of partial solutions can be represented by an OBDD of width at most $3$.
% \end{proposition}

% A tree decomposition is immediate if whenever a node $u$ is an introduce node, introducing a vertex $v$, all edges incident with $v$ in the graph $G_u$ are introduced. 


% \begin{conjecture}
% \label{conjecture:CliquesSmallWidth}
% If $G$ is a clique, and $T$ is an immediate tree decomposition of $G$, then the symbolic algorithm has "Small" width during the whole process.  
% \end{conjecture}

% \begin{conjecture}
% Complete Bipartite graphs? 
% \end{conjecture}


% Now, the challenge is to generalize this to a larger of class of graphs. 

% \begin{conjecture}
% \label{generalization:}
% What happens for split graphs? 
% \end{conjecture}


% \begin{conjecture}
% \label{generalization:}
% What happens for graphs of small neighborhood diversity?
% \end{conjecture}

% \begin{conjecture}
% What happens with graphs that are covered by $k$ bipartite graphs? 
% \end{conjecture}

% \todo[inline]{\url{https://arxiv.org/pdf/1201.3091.pdf}}




